\SourceLeaf{254}%
etiam hic utemur, scilicet ut pro quovis puncto superficiei \(r\) sit distantia minima a puncto determinato, atque \(\varphi\) angulus in hoc puncto inter elementum primum ipsius \(r\) atque directionem fixam. Ita habemus \(E'=1\), \(F'=0\), \(\omega'=90^0\): statuemus insuper \(\BareRoot{G'}=m\), ita ut elementum lineare quodcunque fiat \(=\BareRoot{(\mathrm{d}r^2+mm\mathrm{d}\varphi^2)}\). Hinc quatuor aequationes in art. praec. pro \(\alpha,\beta,\gamma,\delta\) erutae, suppeditant:
\[
\begin{aligned}
\BareRoot{E}\mathord{.}\cos(\omega-\psi)&=\frac{\mathrm{d}r}{\mathrm{d}p}\quad .\ .\ .\ .\ .\ .\ .\ .\ .\ .\ .\ .\quad(1)\\
\BareRoot{G}\mathord{.}\cos\psi&=\frac{\mathrm{d}r}{\mathrm{d}q}\quad .\ .\ .\ .\ .\ .\ .\ .\ .\ .\ .\ .\quad(2)\\
\BareRoot{E}\mathord{.}\sin(\psi-\omega)&=m\mathord{.}\frac{\mathrm{d}\varphi}{\mathrm{d}p}\quad .\ .\ .\ .\ .\ .\ .\ .\ .\ .\ .\quad(3)\\
\BareRoot{G}\mathord{.}\sin\psi&=m\mathord{.}\frac{\mathrm{d}\varphi}{\mathrm{d}q}\quad .\ .\ .\ .\ .\ .\ .\ .\ .\ .\ .\quad(4)
\end{aligned}
\]
Ultima et penultima vero has
\[
\begin{aligned}
EG-FF&=E(\frac{\mathrm{d}r}{\mathrm{d}q})^2-2F\mathord{.}\frac{\mathrm{d}r}{\mathrm{d}p}\mathord{.}\frac{\mathrm{d}r}{\mathrm{d}q}+G(\frac{\mathrm{d}r}{\mathrm{d}p})^2\quad .\ .\ .\ .\ .\ .\quad(5)\\
(E\mathord{.}\frac{\mathrm{d}r}{\mathrm{d}q}-F\mathord{.}\frac{\mathrm{d}r}{\mathrm{d}p})\mathord{.}\frac{\mathrm{d}\varphi}{\mathrm{d}q}&=(F\mathord{.}\frac{\mathrm{d}r}{\mathrm{d}q}-G\mathord{.}\frac{\mathrm{d}r}{\mathrm{d}p})\mathord{.}\frac{\mathrm{d}\varphi}{\mathrm{d}p}\quad .\ .\ .\ .\ .\ .\ .\quad(6)
\end{aligned}
\]

Ex his aequationibus petenda est determinatio quantitatum \(r,\varphi,\psi\) et (si opus videatur) \(m\), per \(p\) et \(q\): scilicet integratio aequationis (5) dabit \(r\), qua inventa integratio aequationis (6) dabit \(\varphi\), atque alterutra aequationum (1), (2) ipsam \(\psi\): denique \(m\) habebitur per alterutram aequationum (3), (4).

Integratio generalis aequationum (5), (6) necessario duas functiones arbitrarias introducere debet, quae quid sibi velint facile intelligemus, si perpendimus, illas aequationes ad casum eum quem hic consideramus non limitari, sed perinde valere, si \(r\) et \(\varphi\) accipiantur in significatione generaliore art. 16, ita ut sit \(r\) longitudo lineae brevissimae ad lineam arbitrariam determinatam normaliter ductae, atque \(\varphi\) functio arbitraria longitudinis eius partis lineae, quae inter lineam brevissimam indefinitam et punctum arbitrarium determinatum intercipitur. Solutio itaque generalis haec omnia indefinite amplecti debet, functionesque arbitrariae tunc demum in definitas abibunt, quando linea illa arbitraria atque functio partium, quam \(\varphi\) exhibere debet, praescriptae sunt. In casu nostro circulus infinite parvus adoptari potest, centrum in eo puncto habens, a quo distantiae \(r\) numerantur, et \(\varphi\) denotabit partes huius circuli ipsas per radium divisas,
